\subsection{The effect on realized transaction prices}
\label{sec:saleprice}

The outcome studied so far is the asking price, which the listing agent sets jointly
with the description and which therefore shares a common author with the treatment.
A natural concern is that part of the estimated association reflects the agent
pricing their own prose rather than the market valuing the property's language. The
cleanest answer is to re-estimate the effect on the price the market actually
cleared. We assemble a panel of recorded sales by collecting sold listings across
the ten metropolitan areas for which a public sold feed carries both the marketing
description and the realized sale price on the same transaction, together with the
sale date and the structural attributes. After restricting to residential
single-unit dwellings with a usable description, a recorded price, and a geocode,
the panel holds roughly two hundred and ninety thousand transactions closed between
2023 and 2026, two orders of magnitude larger than the description samples used
elsewhere in the paper. The treatment is the same pooled leading text direction,
re-estimated on the sold descriptions and held fixed across the estimation, and the
outcome is the log recorded sale price. Because sale prices span three years of
nominal, appreciating-market dollars, we add sale-quarter fixed effects to the
control set so the flexible nuisance learner absorbs the price cycle
\citep{rosen1974hedonic, case1989efficiency}.

The estimand requires care about which attributes to control. The description and
the agent-reported structural fields are authored together, so conditioning on the
listing's own bedroom count or square footage risks blocking a pathway that runs
through the very language we study, a bad control in the sense of
\citet{cinelli2024crash}. We therefore take the exogenous county-assessor record as
the primary control set, which is authored independently of the listing, and read
the resulting estimate as a controlled direct effect of the text net of the
property's structural characteristics. County structural records are available for
eight of the ten markets; for the remaining two, Washington and Atlanta, we report
the agent-reported specification alone. The estimand is further conditional on sale,
since the panel observes only transactions that closed. We do not attempt a
selection correction, for which no credible exclusion restriction exists, and state
the effect as the average text effect among properties that sell, a standard
conditional estimand in this literature \citep{heckman1979selection,
munneke2000empirical}.

The effect survives the change of outcome in every market. Table~\ref{tab:soldprice}
and Figure~\ref{fig:soldprice-forest} report the per-market estimates. On the assessor specification a one-standard-deviation
shift along the leading text direction moves log sale price by between three
and sixteen percent, individually significant in all eight markets under
tract cluster-bootstrap inference, and the heterogeneity mirrors the asking-price
panel: the effect is largest in the older, architecturally heterogeneous markets
(Philadelphia, Chicago, San Francisco, Boston) and smallest in the newer, more
uniform Sun Belt markets (Phoenix, Seattle, Denver, Portland). The sign of the
coefficient is a normalization of the leading component and is not itself
interpretable; the magnitude and its cross-market ordering are. The agent-reported
specification is significant in all ten markets but is uniformly smaller than the
assessor specification in every market where both are estimated, which is the
signature of the bad-control channel we anticipated, since controlling for
text-entangled attributes absorbs part of the effect.

The result is robust along the three dimensions most likely to threaten it.
Spatial dependence is addressed beyond the tract cluster bootstrap by a Conley
spatial HAC standard error and by the self-normalized cluster-$t$ interval of
\citet{ibragimov2010t}, which is the better-calibrated of the available
spatial procedures under cross-fitting \citep{salerno2026spatially}; the effect
excludes zero in all eight assessor markets under every one of these procedures.
Dependence across markets and across calendar quarters is addressed jointly by
pooling the ten markets into a single panel with metro and sale-quarter fixed
effects and reporting the two-way cluster-robust standard error of
\citet{cameron2011robust} by metro and by sale-quarter: the pooled
per-standard-deviation effect is $0.061$ with a two-way-clustered interval of
$[0.026, 0.097]$ that excludes zero. We read this as a robustness check and not as
the headline, because with ten metropolitan clusters the metro dimension sits well
below the count at which two-way cluster-robust inference is dependable
\citep{cameron2015practitioner}; the interval uses a $t$ reference on nine degrees
of freedom, and the metro-clustered standard error alone, which inflates the
influence-function standard error roughly twentyfold, already accounts for almost
the entire two-way width, so the residual quarter dimension is not what carries the
result. The
concern that the estimated text direction is a generated regressor is addressed by
the split-sample discipline of \citet{egami2022texts}: estimating the leading
direction on a random half of the listings and the effect on the disjoint half
returns the same per-market estimates to within a thousandth, so the effect is not
an artifact of discovering and estimating the text measure on the same data.
Finally, the leading text direction is almost uncorrelated with the sale quarter,
with an $R^2$ below two percent in every market, so the effect is not a repackaged
market-cycle trend and the sale-quarter controls are not load-bearing. Because the
sold descriptions arrive truncated at a fixed character length, we report the
assessor estimate both raw and corrected by the validation-sample attenuation factor
of Appendix~\ref{sec:truncation}; the correction widens the effects modestly and
changes no conclusion. That the effect is present, heterogeneous in the same pattern,
and robust on the transaction the market actually cleared, in every market where that
transaction is observable, is the strongest evidence the available data affords that
the listing-text association is a feature of market valuation rather than an artifact
of how agents price their own language.

\begin{table}[H]
\centering
\caption{Effect of the leading listing-text direction on log realized sale price, by metropolitan area. Estimates are ridge partially-linear DML on the single-unit subset with sale-quarter fixed effects; the assessor arm controls for exogenous county-record attributes (the primary, controlled-direct-effect specification) and the listing arm for agent-reported attributes (a bad-control sensitivity bound). Intervals are tract cluster-bootstrap ($B=300$); $\hat\theta_{\text{corr}}$ divides the assessor estimate by the validation-sample truncation factor. $\hat\theta_{\text{listing}}$ and $\hat\theta_{\text{corr}}$ are reported as point estimates, without a separate interval of their own. Estimates are oriented so the effect is positive; all are individually significant.}
\label{tab:soldprice}
\small
\begin{tabular}{lrrrrr}
\toprule
Market & $n$ & $\hat\theta_{\text{assessor}}$ & 95\% CI (cluster) & $\hat\theta_{\text{listing}}$ & $\hat\theta_{\text{corr}}$ \\
\midrule
Philadelphia & $29{,}094$ & $0.159$ & $[0.146, 0.169]$ & $0.140$ & $0.196$ \\
Chicago & $20{,}245$ & $0.148$ & $[0.136, 0.158]$ & $0.102$ & $0.182$ \\
San Francisco & $6{,}585$ & $0.115$ & $[0.098, 0.131]$ & $0.076$ & $0.141$ \\
Boston & $3{,}337$ & $0.111$ & $[0.093, 0.132]$ & $0.057$ & $0.137$ \\
Washington & $9{,}846$ & -- & -- & $0.060$ & -- \\
Atlanta & $23{,}603$ & -- & -- & $0.069$ & -- \\
Seattle & $17{,}363$ & $0.059$ & $[0.053, 0.065]$ & $0.049$ & $0.073$ \\
Denver & $22{,}000$ & $0.060$ & $[0.052, 0.068]$ & $0.043$ & $0.074$ \\
Portland & $21{,}653$ & $0.064$ & $[0.056, 0.069]$ & $0.044$ & $0.080$ \\
Phoenix & $39{,}003$ & $0.029$ & $[0.025, 0.034]$ & $0.020$ & $0.036$ \\
\bottomrule
\end{tabular}
\end{table}


\begin{figure}[H]
\centering
\begin{tikzpicture}
\begin{axis}[
    causalre, width=0.62\textwidth, height=7.4cm,
    ytick={10,9,8,7,6,5,4,3,2,1}, yticklabels={Philadelphia,Chicago,San Francisco,Boston,Portland,Denver,Seattle,Phoenix,Atlanta,Washington},
    xlabel={$\hat\theta$ per $\sigma$ (log sale price)},
    xmin=-0.008, xmax=0.185, xtick={0,0.05,0.10,0.15},
    scaled x ticks=false, xticklabels={$0$,$0.05$,$0.10$,$0.15$},
    enlarge y limits=0.07,
    extra x ticks={0}, extra x tick style={grid=major, grid style={black, dashed, thin}}, extra x tick labels={},
    legend style={at={(0.97,0.06)}, anchor=south east}, legend cell align=left,
]
\draw[thin,gray] (axis cs:0.146,10.17)--(axis cs:0.169,10.17);
\draw[thin,gray] (axis cs:0.136,9.17)--(axis cs:0.158,9.17);
\draw[thin,gray] (axis cs:0.098,8.17)--(axis cs:0.131,8.17);
\draw[thin,gray] (axis cs:0.093,7.17)--(axis cs:0.132,7.17);
\draw[thin,gray] (axis cs:0.056,6.17)--(axis cs:0.069,6.17);
\draw[thin,gray] (axis cs:0.052,5.17)--(axis cs:0.068,5.17);
\draw[thin,gray] (axis cs:0.053,4.17)--(axis cs:0.065,4.17);
\draw[thin,gray] (axis cs:0.025,3.17)--(axis cs:0.034,3.17);
\draw[thin,gray] (axis cs:0.131,9.83)--(axis cs:0.150,9.83);
\draw[thin,gray] (axis cs:0.093,8.83)--(axis cs:0.112,8.83);
\draw[thin,gray] (axis cs:0.066,7.83)--(axis cs:0.087,7.83);
\draw[thin,gray] (axis cs:0.045,6.83)--(axis cs:0.068,6.83);
\draw[thin,gray] (axis cs:0.037,5.83)--(axis cs:0.051,5.83);
\draw[thin,gray] (axis cs:0.036,4.83)--(axis cs:0.049,4.83);
\draw[thin,gray] (axis cs:0.045,3.83)--(axis cs:0.054,3.83);
\draw[thin,gray] (axis cs:0.016,2.83)--(axis cs:0.025,2.83);
\draw[thin,gray] (axis cs:0.061,2.00)--(axis cs:0.078,2.00);
\draw[thin,gray] (axis cs:0.053,1.00)--(axis cs:0.066,1.00);
\addplot[only marks, mark=*, mark size=1.9pt, color=cNYC] coordinates {(0.159,10.17) (0.148,9.17) (0.115,8.17) (0.111,7.17) (0.064,6.17) (0.060,5.17) (0.059,4.17) (0.029,3.17)};
\addlegendentry{assessor controls}
\addplot[only marks, mark=o, mark size=2.1pt, line width=0.8pt, color=cYEL] coordinates {(0.140,9.83) (0.102,8.83) (0.076,7.83) (0.057,6.83) (0.044,5.83) (0.043,4.83) (0.049,3.83) (0.020,2.83) (0.069,2.00) (0.060,1.00)};
\addlegendentry{listing controls}
\end{axis}
\end{tikzpicture}
\caption{Per-market effect of the leading listing-text direction on log realized sale price, oriented so the effect is positive. Filled blue markers use exogenous county-assessor structural controls (the primary, controlled-direct-effect specification); open amber markers use agent-reported listing attributes (the bad-control sensitivity bound). Gray whiskers are $95\%$ tract cluster-bootstrap intervals ($B=300$); the dashed line marks the null. Washington and Atlanta lack a county structural source and carry only the listing specification. The listing effect lies inside the assessor effect in every market where both are estimated, the signature of controlling for text-entangled attributes. Numerical values appear in Table~\ref{tab:soldprice}.}
\label{fig:soldprice-forest}
\end{figure}

